% Minimal Nature-MI-flavored stand-in. Nature uses a custom in-house
% LaTeX template (sn-jnl) that ships with paid submission packages
% and isn't redistributable; this generic-article approximation is
% enough for pandoc to compile paper.md with a Nature-style visual
% rhythm (single column body, large title block, abstract before any
% headings).
\documentclass[11pt,a4paper]{article}
% pdflatex cannot set Chinese, Japanese or Korean text; generation/paper.py
% compiles such a paper with XeLaTeX and passes fi-cjk-font.
\usepackage{iftex}
\ifXeTeX
  \usepackage{fontspec}
$if(fi-cjk-font)$
  \usepackage{xeCJK}
  \setCJKmainfont{$fi-cjk-font$}
  \setCJKsansfont{$fi-cjk-font$}
  \setCJKmonofont{$fi-cjk-font$}
$endif$
\else
  \usepackage[utf8]{inputenc}
\fi
\usepackage[margin=1.1in]{geometry}
\usepackage{graphicx}
% Required for pandoc-rendered markdown tables (pandoc emits \begin{longtable}).
\usepackage{longtable}
% Required for pandoc-rendered table column widths. Pandoc emits column
% specs like ``p{(\linewidth - 6\tabcolsep) * \real{0.2667}}`` whose
% arithmetic only parses with calc loaded. Without this, any markdown
% table with relative column widths fails the LaTeX compile with
% ``Missing number, treated as zero`` at the longtable \begin{}.
\usepackage{calc}
\usepackage{array}
% pandoc emits \def\LTcaptype{none} for caption-less tables; longtable
% then calls \refstepcounter{none} internally, so the counter must exist.
\newcounter{none}
\usepackage{hyperref}
\usepackage{amsmath, amssymb}
\usepackage{booktabs}
\usepackage{microtype}
% Don't pair authblk with titling — both patch \author/\maketitle and
% the combo trips LaTeX's "Missing \begin{document}" on \author[N]{}.
% titling alone is enough for the Nature-style flush-left author block;
% we render the affiliation as the second line of \author{} below.
\usepackage{titling}
\usepackage{titlesec}

% Pandoc-compatibility preamble — see generic/template.tex for the
% full rationale.
\providecommand{\tightlist}{%
  \setlength{\itemsep}{0pt}\setlength{\parskip}{0pt}}
\providecommand{\passthrough}[1]{#1}
% A figure fills the line width (its column, in two columns) but is never
% taller than 45% of the text height, so it shares a page with text rather
% than being moved on to the next page and leaving this one half empty.
\makeatletter
\newsavebox{\FI@figbox}
% With a page limit (generation/paper.py sets fi-tight), a figure is at most
% 33% of the text height, so more text shares its page.
\providecommand{\pandocbounded}[1]{%
  \sbox{\FI@figbox}{#1}%
$if(fi-tight)$
  \Gscale@div\@tempa{0.33\textheight}{\dimexpr\ht\FI@figbox+\dp\FI@figbox\relax}%
$else$
  \Gscale@div\@tempa{0.45\textheight}{\dimexpr\ht\FI@figbox+\dp\FI@figbox\relax}%
$endif$
  \Gscale@div\@tempb{\linewidth}{\wd\FI@figbox}%
  \ifdim\@tempb\p@<\@tempa\p@\let\@tempa\@tempb\fi
  \noindent\scalebox{\@tempa}{\usebox{\FI@figbox}}}
% A figure may also stay where it is written (h), and a page of floats must
% be 80% full, so a figure is never put alone on a page half of it blank.
\def\fps@figure{htbp}
\makeatother
\renewcommand{\topfraction}{0.85}
\renewcommand{\bottomfraction}{0.6}
\renewcommand{\textfraction}{0.1}
\renewcommand{\floatpagefraction}{0.8}
% A figure is a float, so LaTeX may hold it back to a later page. The
% reference list ends the paper's body, and a figure held past it printed
% among the references instead of beside the text that discusses it.
% generation/paper.py puts \FIfloatbarrier just before that list, and any
% figure still waiting is placed before it. With nothing waiting the
% barrier does nothing at all, so a paper whose figures already fit keeps
% its layout unchanged. \@deferlist is LaTeX's own list of floats held for
% a later page; placeins is deliberately not used, because it is absent
% from TeX installs FI otherwise supports.
\makeatletter
\newcommand{\FIfloatbarrier}{\par\ifx\@deferlist\@empty\else\clearpage\fi}
\makeatother
$highlighting-macros$

% Nature's title is larger than the body but not poster-sized; \Large
% bold matches the visual weight of the Nature-MI cover page.
\pretitle{\begin{flushleft}\Large\bfseries}
\posttitle{\par\end{flushleft}\vskip 0.4em}
\preauthor{\begin{flushleft}\normalsize}
\postauthor{\par\end{flushleft}}
\predate{\begin{flushleft}\small}
\postdate{\par\end{flushleft}\vskip 0.6em}

% Nature uses bold section heads, sans-serif feel; titlesec gives us
% a bold black heading without renaming the article-class hierarchy.
\titleformat{\section}{\normalfont\large\bfseries}{\thesection}{0.6em}{}

\hypersetup{pdftitle={$title$}, pdfauthor={$if(fi-author)$$fi-author$$else$Frontier Insight$endif$}$if(keywords)$, pdfkeywords={$for(keywords)$$keywords$$sep$, $endfor$}$endif$}

\title{$title$}
% Author block from the output.author fields (see generic/template.tex).
\author{$if(fi-author)$$fi-author$$else$Frontier Insight$endif$$if(fi-affiliation)$\\{\small $fi-affiliation$}$endif$$if(fi-email)$\\{\small $fi-email$}$endif$$if(fi-url)$\\{\small $fi-url$}$endif$}
\date{$date$}

% A visual check that finds a last page holding only a line or two
% recompiles with fi-extra-lines set: the text area grows that many lines and
% the footer moves up by the same amount, so it stays where it was
% (generation/_visual_check.py).
$if(fi-extra-lines)$
\AtBeginDocument{\addtolength{\textheight}{$fi-extra-lines$\baselineskip}\addtolength{\footskip}{-$fi-extra-lines$\baselineskip}}
$endif$

\begin{document}
\maketitle

$if(abstract)$
\begin{abstract}
$abstract$
\end{abstract}
$endif$
$if(keywords)$
{\small\raggedright\noindent\textbf{Keywords:} $for(keywords)$$keywords$$sep$, $endfor$\par}\medskip
$endif$

$body$

\end{document}
